\begin{figure}[p]\centering
\includegraphics[width=\textwidth]{figures/Fig1A_native_WT.png}\\[1mm]
\includegraphics[width=\textwidth]{figures/Fig1B_native_SCA3.png}\\[1mm]
\includegraphics[width=.48\textwidth]{figures/Fig1C_ISI_endpoints.pdf}\hfill
\includegraphics[width=.48\textwidth]{figures/Fig1D_ISI_reproduction.pdf}
\caption{\textbf{Recorded voltage and historical endpoint descriptors.} (A,B) Native-voltage panels from the channel-validated archive audit for WT\_14 and SCA3\_04 at 400 pA, zero-based ABF sweep 8. Curated peaks and the independently selected command trace are displayed with the original axes. The nonspiking portions are not included in the revised conditional metrics. Panels are layout crops of the complete audit figures without alteration of the plotted signal. These two cells illustrate recordings, not independent group replicates or voltage fits. (C) Historical mean ISI at $q=0.75$ for 12 WT and six SCA3 cells; horizontal bars show medians. (D) Corresponding best-model versus experimental mean ISI. These are archived current-evoked support-window statistics, not spontaneous firing or a common-current comparison.}\label{fig:experiment}
\end{figure}
\begin{figure}[p]\centering
\includegraphics[width=.48\textwidth]{figures/Fig2A_model_traces.pdf}\hfill
\includegraphics[width=.48\textwidth]{figures/Fig2B_state_projection.pdf}\\[2mm]
\includegraphics[width=.48\textwidth]{figures/Fig2C_ISI_cdf.pdf}\hfill
\includegraphics[width=.48\textwidth]{figures/Fig2D_affinities.pdf}
\caption{\textbf{Conditional state and ISI distributions on matched trajectories.} Best-fit WT\_14--SCA3\_04 at $q=0.75$. (A) HR $x(t)$ at $p=0$, 0.50, and 1, with vertical display offsets; units are model units, not mV. (B) Standardized $x$--$z$ projections of active-cycle samples; all three coordinates enter the distance. (C) Empirical log-ISI cumulative distributions from the same accepted intervals. (D) Archived state and ISI affinity curves on the 21-node grid. At $p=0.50$, state affinity is positive while ISI affinity is negative. This selected example explains the construction; cohort conclusions use all retained parameter combinations.}\label{fig:mechanism}
\end{figure}
\begin{figure}[p]\centering
\includegraphics[width=.80\textwidth]{figures/Fig3A_delta_q.pdf}\\[3mm]
\includegraphics[width=.80\textwidth]{figures/Fig3B_complete_pairs.pdf}
\caption{\textbf{Relative-current sensitivity.} (A) Median pair differences for both aggregation procedures. Shaded bands show the IQR across dependent cell pairs, not biological confidence intervals. (B) The same 29 scenario-first pairs evaluable at all six settings, with individual curves and their median. Twenty-six remain negative throughout. Complete-case display does not remove missing pairs from the attempted denominators in the tables. All settings use the same state normalization.}\label{fig:q}
\end{figure}
\begin{figure}[p]\centering
\includegraphics[width=.48\textwidth]{figures/Fig4A_oriented_null.pdf}\hfill
\includegraphics[width=.48\textwidth]{figures/Fig4B_endpoint_contrasts.pdf}\\[2mm]
\includegraphics[width=.48\textwidth]{figures/Fig4C_matching_coverage.pdf}\hfill
\includegraphics[width=.48\textwidth]{figures/Fig4D_fixed_current.pdf}
\caption{\textbf{Controls delimit the ordering.} (A) Fractions of evaluable pairs with negative differences, scenario-first; counts and missing denominators appear in Table~\ref{tab:null}. (B) Pair-level state and log-ISI endpoint contrasts at $q=0.75$. (C) Distinct cross-genotype pairs with at least one matched control and unique matched within-genotype pairs. Both endpoint-distance ratios must be within 1.25; all five matched cross-genotype pairs at $q=0.75$ contain SCA3\_04. (D) Relative-current versus fixed-current differences on exactly the same 23 evaluable pairs. Thirteen change from negative to positive, none in the opposite direction. These controls describe model-path dependence rather than causal genotype effects.}\label{fig:controls}
\end{figure}
